% Faithful transcription — original French. Transcribed from the original first printing:
% Hieronymus Georg Zeuthen, "Note sur la théorie de surfaces réciproques",
% Mathematische Annalen, Band 4 (1871), pp. 633–637.
% Scan: Internet Archive (sim_mathematische-annalen_1871_4).
%
% \origpage uses the printed page numbers (633–637). Running heads and page numbers are omitted.
%
% Cross-page typography and rendering decisions:
% - Summation operator uses capital Greek letter \Sigma (not \sum), following Clebsch/Zeuthen.
% - Singularities notation: \chi, \chi' for close-points (points-clos) and close-planes (plans-clos);
%   \eta, \eta' for binode/biplane singularities; \xi, \xi'; \varrho, \varrho'; \gamma, \gamma'.
% - German quotation marks „..." are preserved as literal Unicode per HOUSESTYLE R22.
% - Printer's errors faithfully reproduced per HOUSESTYLE R4:
%     p. 635: damaged digit in "page 4"; missing accent in "singularites"; missing closing quote.
% See notation.md for the complete specification.
\documentclass{article}
\usepackage{readmasters}
\begin{document}

\origpage{633}
\section*{Note sur la théorie de surfaces réciproques.}

\subsection*{Par H.~G. Zeuthen à Copenhague.}

Dans mes deux mémoires au commencement de ce volume des „Annalen“ j'ai attribué aux surfaces dont je m'occupe certaines singularités que j'ai énumérées dans l'introduction du premier de ces mémoires. J'ai supposé qu'aucune autre singularité ne s'y ajoutât, ou bien que les singularités se présentassent de la manière la plus générale. Or, comme cette manière la plus générale est différente suivant qu'on regarde la surface comme lieu de points ou comme enveloppe de plans, il faut diviser les singularités en \emph{singularités de définition ponctuelle} et en \emph{singularités de définition tangentielle}. Les singularités de définition ponctuelle sont celles qui satisferont de la manière la plus générale à leur définition si l'on regarde la surface comme lieu de points, et les singularités de définition tangentielle sont celles qui y satisferont de la manière la plus générale si l'on regarde la surface comme enveloppe de plans.

Nous n'avons pas dit expressément quelles définitions de singularités sont ponctuelles et quelles tangentielles, sans quoi les définitions ne sont pas complètes. Nous supléerons ici à cette manque en montrant que la division que nous avons suivie, sans le dire, est la plus naturelle. Comme une surface regardée comme lieu de points n'a pas, en général, des courbes doubles et cuspidales, mais bien des séries de plans tangents doubles et stationnaires, les définitions des singularités ordinaires des plans tangents doubles et stationnaires ($b', k', t', \gamma', q', \varrho', c', h', \beta', r', \sigma'$) doivent être ponctuelles, et celles des singularités ordinaires de la courbe double et de la courbe cuspidale ($b, k, t, \gamma, q, \varrho, c, h, \beta, r, \sigma$) doivent être tangentielles.${}^{*)}$ Une

\textbf{*)} Alors on sait que les $\beta$ points cuspidaux de la courbe cuspidale se trouvent sur la courbe double, et que les $\gamma$ points cuspidaux de la courbe double se trouvent sur la courbe cuspidale, et ces points sont donc les mêmes que ceux dont MM.~Salmon et Cayley désignent les nombres de la même manière. M.~Cayley m'a fait observer que les deux courbes peuvent avoir d'autres points cuspidaux. Ces nouveaux points singuliers demanderaient une nouvelle étude. --- J'aurais dû dire expressément que ma définition de $k$ (nombre Plückerien des génératrices doubles d'un cône projetant la courbe double) est un peu différente de celle des MM.~Salmon et Cayley, ce qui explique une discordance apparente entre leurs formules et les miennes. J'écris $k - f - 3t - \Sigma(\cdots)$ où M.~Cayley écrit seulement $k$.

\origpage{634}
courbe double, une fois attribuée à une surface regardée comme lieu de points, sera en général douée de points-pinces, pendant que la courbe double d'une surface regardée comme enveloppe de plans n'en a pas généralement. Pour cette raison on doit donner à ces $j$ singularités une définition ponctuelle, et aux $j'$ plans-pinces une définition tangentielle. De même les $\chi$ points-clos sont de définition ponctuelle, et les $\chi'$ plans-clos de définition tangentielle. Nous avons donné expressément aux points coniques et aux singularités qui y appartiennent une définition ponctuelle, et aux plans qui sont tangents le long d'une courbe et aux singularités qui y appartiennent une définition tangentielle. C'est en regardant les définitions de $f$ et de $i$ comme ponctuelles, et celles de $f'$ et de $i'$ comme tangentielles, qu'on a les équations $f = f'$ et $i = i'$.

Nous faisons encore une autre distinction: les propriétés ponctuelles d'une singularité sont celles qui ont rapport aux points de la surface, et les propriétés tangentielles celles qui ont rapport à ses plans tangents. $t'$ indique, par exemple, le nombre des plans qui ont la propriété ponctuelle de rencontrer la surface en une courbe douée de trois points doubles, et la propriété tangentielle d'être plans tangents triples, soit de la surface elle-même, soit de l'enveloppe des plans tangents doubles. Les propriétés ponctuelles des singularités de définition ponctuelle, ainsi que les propriétés tangentielles des singularités de définition tangentielle, sont indiquées dans les définitions elles-mêmes: mais il est souvent très-difficile de trouver les propriétés ponctuelles des singularités de définition tangentielle, et les propriétés tangentielles des singularités de définition ponctuelle, et ces propriétés peuvent être assez compliquées. Pour connaître suffisamment les propriétés tangentielles d'un point-pince, par exemple, il faut avoir étudié la singularité de l'enveloppe des plans tangents stationnaires causée par le point-pince. Pour la trouver on a besoin d'étudier les singularités que l'Hessienne et sa courbe d'intersection avec la surface donnée ont à ce même point.

Si l'on regarde, dans une recherche, une surface comme lieu de ses points on doit connaître les propriétés ponctuelles des singularités qu'on lui attribue, et si l'on regarde une surface comme enveloppe de ses plans tangents on doit connaître les propriétés tangentielles de ses singularités. Dans une théorie des surfaces réciproques on a besoin de connaître, à la fois, les propriétés ponctuelles et les propriétés tangentielles des singularités qu'on attribue aux surfaces. N'ayant fait une étude assez profonde ni des propriétés tangentielles des singularités $j, \chi, \eta$ et $\xi$, dont la définition est ponctuelle, ni des propriétés ponctuelles des singularités $j', \chi', \eta'$ et $\xi'$, dont la définition est tangentielle, je dois prier les lecteurs de mes mémoires \emph{d'ajouter au}

\origpage{635}
mémoire „Sur deux surfaces dont les points se correspondent un-à-un“, où les surfaces sont regardées comme lieux de points, la restriction suivante:

„\emph{Les surfaces $[F_1]$ et $[F_2]$ sont dénuées des singularités $j', \chi', \eta'$ et $\xi'$“},

et à l'introduction du mémoire „Sur les droites multiples des surfaces“ la restriction suivante:

„\emph{Les équations à la page 4\ednote{In the printed text the digit 4 has a heavy ink defect or batter; the context of the memoir establishes page 4.} sont justes si la surface est dénuée des singularites\ednote{Printed singularites without accent; an apparent misprint for singularités.} $j', \chi', \eta'$ et $\xi'$, et les équations réciproques si elle est dénuée des singularités $j, \chi, \eta$ et $\xi$. Elles seront donc applicables à la fois à une surface dénuée de toutes ces singularités.}\ednote{The quotation opened with „Les équations lacks a closing quotation mark in the print.}

Ni à l'un ni à l'autre de ces endroits les restrictions n'altéreront les exposés des résultats; mais il est possible que les singularités dont je parle, si les surfaces en avaient, introduiraient de nouveaux termes dans les équations. Quant aux recherches qui font l'objet principal de mon premier mémoire, je n'ai pas besoin d'y ajouter aussi une restriction; car les singularités $j$ et $j'$ ne s'y présentent que d'une manière spéciale.

Je ne dis pas que les formules à la page 4 cessent d'être vraies au delà des limites indiquées ici, mais seulement que \emph{ma déduction} de ces formules (que je n'ai pas exposée mais où je fais usage, soit des procédés de M.~Salmon, soit du principe de correspondance) est alors insuffisante. Pour les singularités $j$ et $j'$, $\chi$ et $\chi'$ je ne fais que reproduire les formules de M.~Cayley (On reciprocal Surfaces. Phil.\ Trans.\ 1869, p.~202), et quant aux singularités $\eta$ et $\eta'$ mes résultats sont d'accord avec ceux que ce savant a trouvé par rapport aux „binodes“ et aux „biplans“; mais je ne sais pas si M.~Cayley veut donner à ses définitions les mêmes degrés de généralité que j'ai indiqués ici, ou s'il se borne aux cas où les singularités de définition tangentielle n'ont pas d'autres propriétés ponctuelles que celles qui appartiennent déjà aux singularités ponctuelles qu'il attribue aux surfaces, et réciproquement. Cela a lieu pour les surfaces cubiques qui font un objet principal de ses recherches; mais déjà à l'étude d'une surface de $4^{\text{me}}$ ordre douée d'une conique double et à laquelle on donne une définition ponctuelle on rencontre des points-pinces qui ont des propriétés tangentielles plus compliquées. Les nombres des singularités de cette surface, tirées des équations dont je viens de parler (Cayley p.~202, Zeuthen p.~4), ne satisfont pas à l'équation que donne M.~Cayley à la page 228. Nous montrerons dans ce qui suit un autre fait qui sera en discordance avec cette dernière équation (les autres supposées justes) si l'on attribue à la surface des points-clos définis comme je l'ai fait ici.

\origpage{636}
En bornant l'étude des surfaces réciproques aux surfaces dénuées des singularités $j, j', \chi, \chi', \eta, \eta', \xi, \xi'$ il sera moins difficile de déterminer tous les coefficients de la $26^{\text{me}}$ équation${}^{*)}$ qui, à côté des 3 équations (I), des 11 équations (II)---(IV) et (VIII)---(IX) et des 11 équations réciproques, ont lieu entre les singularités d'une surface. En effet, cette $26^{\text{me}}$ équation sera l'équation (13) [ou (14)] à la page 46 de mon deuxième mémoire.

Je donnerai ici une autre déduction de la même équation où je ne ferai pas usage de la théorie de deux surfaces dont les points se correspondent un-à-un, mais seulement du principe de correspondance. Nous appliquerons ce principe à \emph{trouver le nombre des plans qui sont à la fois tangents à la surface et osculateurs à la courbe cuspidale à un même point de celle-ci}. Nous prenons alors pour les points correspondants ($x$ et $u$) dont on cherche les coïncidences les points où le plan tangent et le plan osculateur à un même point de la courbe cuspidale rencontrent une droite fixe. A un point $x$ correspondent $\sigma$ points $u$, et, si l'on désigne par $m$ la classe de la développable (d'ordre $r$) enveloppe des plans osculateurs de la courbe cuspidale, à un point $u$ correspondent $m$ points $x$. Il y a donc $\sigma + m$ coïncidences d'un point $x$ avec un point correspondant $u$. Ces coïncidences ont lieu $1^0$ aux $r$ points où la droite fixe rencontre des tangentes de la courbe cuspidale, et $2^0$ aux points où elle rencontre les plans cherchés. Le nombre de ces plans est donc $\sigma + m - r$.

Une partie des plans trouvés, disons préalablement $\Sigma'$, coïncident avec les plans qui sont tangents le long d'une courbe; les points de contact des autres auront la propriété réciproque à celle des plans: celle d'être à la fois points de contact de plans tangents stationnaires avec la surface et points de contact des mêmes plans avec l'arête de rebroussement de l'enveloppe des plans tangents stationnaires. En déterminant le nombre de ces points d'une manière analogue on trouve l'équation suivante
\[
\sigma + m - r - \Sigma' = \sigma' + m' - r' - \Sigma,
\]
où $m'$ est l'ordre de l'arête de rebroussement de l'enveloppe des plans tangents stationnaires, $\Sigma$ le nombre des points qui ont la propriété indiquée et qui se trouvent aux points coniques de la surface. Or on a d'après les relations de Plücker
\[
m = \beta + 3r - 3c \quad , \quad m' = \beta' + 3r' - 3c',
\]
et $\Sigma$ sera égal à la somme des nombres respectifs des droites qui ont

\textbf{*)} Après avoir exposé les autres 25 équations aux pages 3 et 4 je parle à la page 4 (et à la page 46) de \emph{deux} nouvelles équations (servant à déterminer $\beta$ et $\beta'$); mais ces équations se réduisent, au moyen des 25 autres, à une seule. (Voir le mémoire de M.~Cayley p.~228.)

\origpage{637}
aux points coniques ($\mu$-tuples) un contact $(\mu + 2)$-ponctuel avec la surface (voir à la page 39), ou bien à
\[
\Sigma [\mu(\mu + 1) - 2y - 3z] = \Sigma [2\mu + v].
\]
On aura de même à substituer à $\Sigma'$ le nombre $\Sigma' [2\mu' + v']$. On trouve ainsi
\[
\sigma + 2r - 3c + \Sigma(2\mu + v) = \sigma' + 2r' - 3c' + \Sigma'(2\mu' + v'),
\]
qui est une forme de l'équation cherchée. On peut au moyen des équations à la page 4 la réduire à l'équation (13) à la page 46.

On pourrait appliquer le même procédé à une surface douée des singularités $j, j'$ etc.; mais ne connaissant pas suffisamment les propriétés de ces singularités, comme je l'ai dit, je ne pourrais alors déterminer les coefficients. Toutefois cette recherche peut servir à confirmer que l'introduction de ces singularités dans la théorie des surfaces réciproques a été prématurée. Supposons par exemple que la surface soit douée de $\chi$ points-clos et de $\chi'$ plans-clos. Alors on doit remplacer la formule trouvée par la suivante:
\[
\begin{aligned}
\sigma + 2r - 3c + \Sigma(2\mu + v) + A \cdot \chi \\
= \sigma' + 2r' - 3c' + \Sigma'(2\mu' + v') + A \cdot \chi',
\end{aligned}
\]
où $A$ est un coefficient inconnu mais \emph{entier}; mais l'équation de M.~Cayley à la page 228, combinée avec ses autres formules (à la page 202, mes formules de la page 4) donnerait à $\chi$ et $\chi'$ le coefficient $\frac{25}{3}$.${}^{*)}$

\textbf{*)} Je saisis l'occasion pour répondre à une remarque de M.~Nöther (Nachrichten der Königl.\ Gesellschaft der Wissenschaften zu Göttingen 1871, p.~277): Dans les remarques que j'ai faites à la page 324 du $3^{\text{me}}$ vol.\ des „Mathematische Annalen“ sur la représentation d'une surface sur un plan double je fais, sans le dire je l'avoue, la même supposition que je fais expressément dans mon mémoire sur „deux surfaces dont les points se correspondent un-à-un“, celle qu'aucun des points fondamentaux ne soit doué d'une direction fondamentale.
\end{document}
